\documentclass[10pt,a4paper]{article}

% ---------- Page and font setup ----------
\usepackage[a4paper,top=1.20cm,bottom=1.20cm,left=1.55cm,right=1.55cm]{geometry}
\usepackage[T1]{fontenc}
\usepackage[utf8]{inputenc}
\usepackage[scaled=0.96]{helvet}
\renewcommand{\familydefault}{\sfdefault}
\usepackage{microtype}
\usepackage{graphicx}
\usepackage{tabularx}
\usepackage{enumitem}
\usepackage{xcolor}
\usepackage{titlesec}
\usepackage[hidelinks]{hyperref}
\usepackage{ragged2e}
\usepackage{fancyhdr}
\usepackage{lastpage}

% ---------- PDF metadata ----------
\hypersetup{
  pdftitle={REPLACENAME - Wind Energy Thesis CV},
  pdfauthor={REPLACENAME},
  pdfsubject={Finite element analysis, structural dynamics and wind turbine load modelling}
}

% ---------- Style ----------
\definecolor{heading}{HTML}{202020}
\definecolor{muted}{HTML}{555555}
\definecolor{rulecolor}{HTML}{777777}

\pagestyle{fancy}
\fancyhf{}
\renewcommand{\headrulewidth}{0pt}
\renewcommand{\footrulewidth}{0pt}
\fancyfoot[L]{\scriptsize\color{muted}REPLACENAME}
\fancyfoot[R]{\scriptsize\color{muted}\thepage/\pageref{LastPage}}

\setlength{\parindent}{0pt}
\setlength{\parskip}{0pt}
\setlist[itemize]{leftmargin=1.18em,itemsep=2.5pt,topsep=3pt,parsep=0pt,partopsep=0pt}
\urlstyle{same}
\RaggedRight

\titleformat{\section}
  {\large\bfseries\color{heading}}
  {}{0pt}{}
  [\vspace{-0.22em}\color{rulecolor}\titlerule]
\titlespacing*{\section}{0pt}{0.70em}{0.38em}

\newcommand{\cvjob}[4]{%
  \begin{tabularx}{\linewidth}{@{}Xr@{}}
    \textbf{#1} & \textbf{#2}\\[1pt]
    \textit{#3} & \textit{#4}
  \end{tabularx}
  \vspace{0.05em}
}

\newcommand{\cvproject}[2]{%
  \begin{tabularx}{\linewidth}{@{}Xr@{}}
    \textbf{#1} & \textbf{#2}
  \end{tabularx}
  \vspace{0.08em}
}

\newcommand{\cvedu}[4]{%
  \begin{tabularx}{\linewidth}{@{}Xr@{}}
    \textbf{#1} & \textbf{#2}\\[1pt]
    \textit{#3} & \textit{#4}
  \end{tabularx}
  \vspace{0.03em}
}

\hyphenpenalty=10000
\exhyphenpenalty=10000
\emergencystretch=2em
\renewcommand{\arraystretch}{1.1}
\hyphenation{SOLIDWORKS OpenFAST}

\begin{document}
\fontsize{11.1}{12.7}\selectfont

% ---------- Header ----------
\begin{minipage}[t]{0.79\linewidth}
  \vspace{0pt}
  {\fontsize{21.5}{23.5}\selectfont\bfseries REPLACENAME}\\[2pt]
  {\fontsize{11.55}{13.2}\selectfont\bfseries M.Sc. Wind Energy Engineering | Mechanical Engineer}\\[4pt]
  REPLACELOCATION \enspace|\enspace
  \href{tel:REPLACEPHONE}{REPLACEPHONE} \enspace|\enspace
  \href{mailto:REPLACEEMAIL}{REPLACEEMAIL}\\[1.5pt]
  \href{REPLACELINKEDIN}{REPLACELINKEDIN} \enspace|\enspace
  \href{REPLACEPORTFOLIO}{REPLACEPORTFOLIO}\\[1.5pt]
\end{minipage}
\hfill
\begin{minipage}[t]{0.17\linewidth}
  \vspace{0pt}
  \raggedleft
  \IfFileExists{candidate_photo.jpeg}{%
    \includegraphics[width=2.35cm,height=3.00cm,keepaspectratio]{candidate_photo.jpeg}%
  }{%
    \rule{0pt}{3.00cm}\hspace{2.35cm}%
  }
\end{minipage}

\section{Professional Profile}
M.Sc. Wind Energy Engineering student and mechanical engineer with nearly six years of professional experience in finite-element-based structural and dynamic analysis, mechanical design, manufacturing support and technical documentation. Professional CAE experience includes static structural, modal and harmonic-response analysis in ANSYS. Academic wind-energy work includes aeroelastic simulation of a 20 MW offshore wind turbine with OpenFAST using DNV-ST-0437 and IEC 61400-1 design load cases, supported by Python/MATLAB data processing and automation. Seeking a Master’s thesis at the intersection of wind-turbine structural mechanics, finite-element modelling and component/load assessment.

\section{Selected Academic Wind Energy Projects}
\cvproject{20 MW Offshore Wind Turbine - Structural Load \& Fatigue Analysis}{Academic project}
\begin{itemize}
  \item Conducted OpenFAST aeroelastic simulations for a monopile-supported 20 MW offshore wind turbine based on DNV-ST-0437, including IEC 61400-1 design load cases and evaluation of key structural load channels.
  \item Evaluated turbine response, structural loads and fatigue-relevant channels across environmental and operating conditions using Python- and MATLAB-supported post-processing workflows.
\end{itemize}
\vspace{1pt}

\cvproject
  {\href{REPLACEPROJECTLINK}
  {OpenFAST Workflow Automation \& Output Monitoring}}
  {Academic project}
\begin{itemize}
  \item Developed Python utilities for OpenFAST simulation automation, batch execution and post-processing of simulation outputs and technical files.
  \item Built a Python tool for real-time and post-processing visualization of OpenFAST outputs, supporting efficient review of turbine operating parameters, loads and dynamic response.
\end{itemize}

\cvproject{Wind-Farm Planning and Environmental Assessment}{Academic project}
\begin{itemize}
  \item Completed a site-based wind-farm planning study using windPRO, covering turbine layout, energy-yield estimation, noise assessment and shadow-flicker analysis.
  \item Evaluated bat-curtailment strategies and quantified their influence on expected annual energy production.
\end{itemize}


\section{Professional Experience}
\cvjob{Development Engineer - Solar Tracking Systems \& Automated Mechanisms}{01/2022--02/2023}{Paydar | Solar-Energy Systems}{Tehran}
\begin{itemize}
  \item Designed photovoltaic tracking structures, mounting components and mechanical assemblies in SOLIDWORKS, supported by ANSYS structural analysis and coordination with manufacturing teams.
  \item Developed an electric-motor-driven panel-cleaning mechanism and supported detailed design, fabrication, assembly and successful on-site installation.
  \item Improved manufacturability of mechanical components, contributing to an approximately 10\% reduction in production scrap and raw-material waste; maintained drawings, bills of materials and technical documentation.
\end{itemize}
\vspace{1pt}

\cvjob{Mechanical Engineer - Machinery Development \& Structural Dynamics}{03/2017--12/2021}{TehranRigzar | Mining \& Mineral-Processing Equipment}{Tehran}
\begin{itemize}
  \item Performed static structural, modal and harmonic-response analyses in ANSYS to assess stresses, deformation, resonance and dynamic response of mechanical equipment and structural assemblies.
  \item Designed crushers, conveyors, high-frequency vibrating screens and washing systems; produced detailed fabrication drawings and mechanical documentation in CATIA V5.
  \item Co-developed a high-frequency vibrating screen based on laboratory screening studies, translating operating requirements into structural and mechanical design decisions.
\end{itemize}
\vspace{1pt}

\section{Education}
\cvedu
  {M.Sc. Wind Energy Engineering - Current grade average: 1.9}
  {03/2023--present}
  {Flensburg University of Applied Sciences, Flensburg, Germany}
  {Coursework completed}

\vspace{1.0mm}

\cvedu
  {M.Sc. Mechanical Engineering}
  {09/2014--01/2017}
  {Hakim Sabzevari University, Sabzevar}
  {Control systems and robotics}

\vspace{0.5mm}
\textit{Thesis:} Developed and implemented a MATLAB/Simulink motion-control system for precise movement and positioning of an autonomous underwater welding robot.

\vspace{1.0mm}

\cvedu
  {B.Sc. Mechanical Engineering}
  {09/2008--06/2013}
  {Islamic Azad University, Tehran}
  {Mechanical design and vibration analysis}

\section{Technical Skills}
\begin{tabularx}{\linewidth}{@{}>{\raggedright\arraybackslash}p{4cm}X@{}}
\textbf{FEA \& structural mechanics} & ANSYS Workbench (static structural, modal and harmonic-response analysis), Abaqus, structural dynamics, vibration analysis, stress/deformation assessment\\[2pt]
\textbf{Wind turbine simulation} & OpenFAST, DNV-ST-0437, IEC 61400-1 design load cases, structural-load and fatigue-channel evaluation\\[2pt]
\textbf{Programming \& data} & Python, MATLAB/Simulink, simulation automation, data processing, visualization, technical post-processing\\[2pt]
\textbf{CAD \& engineering} & CATIA V5, SOLIDWORKS, Siemens NX, mechanical design, fabrication drawings, technical documentation\\[2pt]
\textbf{Additional wind tools} & windPRO, QGIS, wind-farm planning, energy-yield estimation, noise and shadow-flicker assessment
\end{tabularx}

\section{Training \& Certifications}
IEEE ANSYS Training \enspace|\enspace HarvardX Python Programming \enspace|\enspace SOLIDWORKS Training \enspace|\enspace CATIA V5 Training

\section{Languages}
\begin{tabularx}{\linewidth}{@{}>{\raggedright\arraybackslash}p{4cm}X@{}}
\textbf{English} & C1; IELTS overall score 7.0\\
\textbf{German} & B1; preparing for certification
\end{tabularx}

\end{document}
